% F3 · Qué se ve y qué no en un depósito con metadatos públicos (D13·02)
\documentclass[border=10pt]{standalone}
% ------------------------------------------------------------------
% Preámbulo común de las figuras del curso C1
% Paleta y tipografía alineadas con ../../estilo.tex
% ------------------------------------------------------------------
\usepackage[utf8]{inputenc}
\usepackage[T1]{fontenc}
\usepackage[spanish,es-noquoting]{babel}
\usepackage{helvet}
\renewcommand{\familydefault}{\sfdefault}
\usepackage{xcolor}
\usepackage{tikz}
\usetikzlibrary{arrows.meta,positioning,decorations.pathreplacing,calc,shapes.geometric,fit,backgrounds}

\definecolor{cHead}{RGB}{68,84,106}
\definecolor{cAcc}{RGB}{0,102,153}
\definecolor{cGris}{gray}{0.42}
\definecolor{cSuave}{RGB}{240,244,248}
\definecolor{cBien}{RGB}{34,124,68}
\definecolor{cMal}{RGB}{176,42,42}
\definecolor{cBienSuave}{RGB}{230,243,234}
\definecolor{cMalSuave}{RGB}{250,235,233}
\definecolor{cAccSuave}{RGB}{225,238,246}

\tikzset{
  caja/.style     = {rectangle, rounded corners=2.5pt, draw=cHead, line width=.7pt,
                     fill=cSuave, align=center, inner sep=6pt, font=\small},
  cajaAcc/.style  = {caja, draw=cAcc, fill=cAccSuave},
  cajaOk/.style   = {caja, draw=cBien, fill=cBienSuave},
  cajaNo/.style   = {caja, draw=cMal, fill=cMalSuave, dashed},
  cajaTenue/.style= {caja, draw=cGris, fill=white},
  rotulo/.style   = {font=\scriptsize\itshape, text=cGris, align=center},
  etiq/.style     = {font=\scriptsize, text=cHead, align=center,
                     fill=white, inner sep=1.5pt},
  fl/.style       = {-{Stealth[length=4.5pt]}, draw=cHead, line width=.7pt},
  flAcc/.style    = {fl, draw=cAcc},
  flNo/.style     = {-{Stealth[length=4.5pt]}, draw=cMal, line width=.7pt, dashed},
  llave/.style    = {decorate, decoration={brace, amplitude=4pt, raise=2pt},
                     draw=cGris, line width=.6pt},
  titulo/.style   = {font=\small\bfseries, text=cHead},
}

% Cabecera de figura: #1 número, #2 título
\newcommand{\figtit}[2]{\node[font=\footnotesize\bfseries, text=cGris, anchor=west]}

\begin{document}
\begin{tikzpicture}

% ---------------- el registro público -------------------------------
\node[cajaOk, text width=6.1cm, align=left, inner sep=9pt] (reg) at (0,0) {
  {\small\bfseries El registro · público}\\[4pt]
  \scriptsize
  Identificador permanente (DOI)\\
  Título y descripción\\
  Población, periodo y criterios de inclusión\\
  Lista de variables y diccionario\\
  Condiciones de acceso y licencia\\
  Procedimiento para solicitarlo\\[3pt]
  \textcolor{cBien}{\bfseries Indexable · citable · encontrable}};

% ---------------- los ficheros --------------------------------------
\node[cajaNo, text width=4.4cm, align=left, inner sep=9pt] (fich) at (8.6,0) {
  {\small\bfseries Los ficheros · restringidos}\\[4pt]
  \scriptsize
  Registros individuales\\
  Datos personales de\\ categoría especial\\[3pt]
  \textcolor{cMal}{\bfseries No se descargan}};

\node[rotulo] at ($(reg.north)!0.5!(fich.north)+(0,.62)$) {un mismo depósito, un mismo identificador};
\draw[cGris, line width=.6pt, dash pattern=on 2pt off 2pt]
  ($(reg.north east)+(.55,.35)$) -- ($(reg.south east)+(.55,-.35)$);

% ---------------- quién llega a cada cosa ---------------------------
\node[cajaTenue, text width=2.5cm] (todos) at (0,-3.1) {Cualquiera\\ en el mundo};
\draw[fl, draw=cBien] (todos) -- (reg.south -| todos)
  node[etiq, pos=.55, text=cBien] {acceso libre};

\node[cajaTenue, text width=2.6cm] (solic) at (8.6,-4.6) {Investigador\\ solicitante};
\node[cajaAcc, text width=4.4cm, align=left, inner sep=7pt] (com) at (8.6,-2.55) {
  {\scriptsize\bfseries Evaluación de la solicitud}\\[2pt]
  \scriptsize órgano que decide\\ criterios publicados\\ acuerdo de uso · plazo};
\draw[flAcc] (solic) -- (com);
\draw[flAcc] (com) -- (fich);

% ---------------- el mensaje ----------------------------------------
\node[rotulo, text width=12cm, anchor=north] at (4.3,-5.9)
  {El catálogo es público; el dato, no. Así se cumplen \textbf{F}, \textbf{I} y \textbf{R} por completo,
   y la \textbf{A} en su versión de acceso controlado, sin publicar un solo registro.};

\end{tikzpicture}
\end{document}
